El proyecto tiene como principal alcance el desarrollo de un sistema de apoyo a la decisión orientado a la predicción de la demanda de medicamentos críticos en una farmacia hospitalaria, esto utilizando los datos históricos de consumo como principal fuente de información. El sistema permitirá estimar la cantidad mensual requerida de medicamentos mediante la implementación de modelos de predicción de demanda, integrando funcionalidades como visualización de resultados, generación de reportes y apoyo en la toma de decisiones relacionadas con la planificación del abastecimiento. Asimismo, el sistema proporcionara al usuario cálculos como el punto de reposición, con el fin de contribuir a una mejor organización y control de los recursos disponibles en los entornos hospitalario.

El desarrollo del proyecto presenta algunas restricciones que deben ser consideradas. En primer lugar, el sistema se basará en un conjunto de datos históricos reales de la Clínica Nuestra señora de los remedios limitado al año 2021, lo que puede afectar la capacidad del modelo para capturar comportamientos actuales o cambios en la demanda a lo largo del tiempo. Además, no se contará con integración en tiempo real con sistemas hospitalarios, por lo que las predicciones se realizarán sobre datos estáticos.